% Source: https://artofproblemsolving.com/community/c3h1869186_a_question_about_symmetry_of_3d_shapes
% Stable identity: aops-topic-1869186-q01
(a) A symmetry maps the cube onto itself, so it fixes the centre $O$ of the cube, and every axis passes through $O$. The axis leaves the cube at a point $P$ of its surface that the rotation fixes. The rotation permutes the faces, the edges and the vertices, and it must map the face, edge or vertex containing $P$ to itself. This gives three types of axis.
\begin{itemize}
\item $P$ is inside a face. That face is mapped to itself and turned about its centre, so $P$ is the centre of the face and the axis joins the centres of two opposite faces. A square is mapped to itself by turns of $90^\circ$: \textbf{order $4$}. There are $\dfrac{6}{2}=3$ such axes.
\item $P$ is inside an edge. The edge is mapped to itself with its ends swapped, so $P$ is its midpoint, and the axis joins the midpoints of two opposite edges. Only the half-turn works: \textbf{order $2$}. There are $\dfrac{12}{2}=6$ such axes.
\item $P$ is a vertex. The axis is a long diagonal of the cube. The three edges at that vertex can be cycled by turns of $120^\circ$: \textbf{order $3$}. There are $\dfrac{8}{2}=4$ such axes.
\end{itemize}

\begin{asy}
usepackage("amsmath");
size(6cm);
defaultpen(fontsize(10pt));
pair P(real x,real y,real z){return (x+0.42*y,z+0.32*y);}
// hidden edges
draw(P(0,1,0)--P(1,1,0),dashed+gray);draw(P(0,1,0)--P(0,0,0),dashed+gray);draw(P(0,1,0)--P(0,1,1),dashed+gray);
// visible edges
draw(P(0,0,0)--P(1,0,0)--P(1,0,1)--P(0,0,1)--cycle,linewidth(0.9));
draw(P(1,0,0)--P(1,1,0)--P(1,1,1)--P(1,0,1),linewidth(0.9));
draw(P(0,0,1)--P(0,1,1)--P(1,1,1),linewidth(0.9));
// face axis (order 4)
draw(P(0.5,0.5,-0.45)--P(0.5,0.5,1.45),red+linewidth(1));
dot(P(0.5,0.5,1),red);dot(P(0.5,0.5,0),red);
label("order $4$",P(0.5,0.5,1.45),N,red);
// vertex axis (order 3)
draw(P(-0.3,-0.3,-0.3)--P(1.3,1.3,1.3),blue+linewidth(1));
label("order $3$",P(1.3,1.3,1.3),NE,blue);
// edge axis (order 2): midpoints of opposite edges (0,0,1)-(1,0,1) and (0,1,0)-(1,1,0)
draw(P(0.5,-0.35,1.35)--P(0.5,1.35,-0.35),deepgreen+linewidth(1));
dot(P(0.5,0,1),deepgreen);dot(P(0.5,1,0),deepgreen);
label("order $2$",P(0.5,1.35,-0.35),SE,deepgreen);
label("\mbox{}",(min(currentpicture,user=true)+max(currentpicture,user=true))/2,black); // keeps the text after the figure black when LaTeX typesets the labels inline
\end{asy}

So the cube has $3+6+4=\boxed{13}$ axes of rotational symmetry: $3$ of order $4$, $6$ of order $2$ and $4$ of order $3$.

(b) Each axis of order $k$ gives $k-1$ non-trivial rotations, and different axes give different rotations. Together with the identity there are
\[
1+3\cdot3+6\cdot1+4\cdot2=\boxed{24}
\]
rotations. Check: a chosen face can be moved to any of the $6$ face positions and then turned in $4$ ways, giving $6\cdot4=24$.

The same method works for other solids. For each kind of point where an axis can leave the solid (centre of a face, midpoint of an edge, a vertex), the order is the number of ways to turn that face, edge or vertex onto itself.
